\section{Introduction}

Heterogeneous graphs are widely used to model real-world systems that contain multiple types of entities and relations, such as academic networks, recommendation platforms, biomedical networks, and social media systems. Compared with homogeneous graphs, heterogeneous graphs provide richer semantic information through relation types and meta-paths. This semantic structure enables heterogeneous graph neural networks (HGNNs), such as HAN, HGT, MAGNN, and HeCo, to learn expressive node representations by aggregating information from multiple semantic neighborhoods \cite{wang2019han,hu2020hgt,fu2020magnn,wang2021heco}. These models have achieved strong performance on node classification and representation learning tasks.

Despite their effectiveness, HGNNs remain vulnerable to adversarial structural perturbations. A small number of malicious edge insertions or deletions can change not only local neighborhoods, but also the high-order semantic neighborhoods induced by meta-paths. This issue is particularly serious in heterogeneous graphs because a perturbation on one relation type may be amplified after meta-path composition and may affect many target-type nodes. As a result, semantic aggregation may propagate misleading information, and attention mechanisms may assign large weights to corrupted semantic views.

Existing robust graph learning methods mitigate structural attacks through edge pruning, probabilistic modeling, graph structure learning, or similarity-based message passing \cite{wu2019gcnjaccard,zhu2019rgcn,jin2020prognn,zhang2020gnnguard,jin2021simpgcn}. Recent robust heterogeneous graph methods further consider typed relations and semantic structures, for example by estimating edge confidence or suppressing suspicious meta-path-induced connections \cite{rohe,hseco,heteguard}. However, most of these methods still construct candidate neighborhoods from the observed heterogeneous topology. When the topology is heavily polluted, topology-only defenses may have limited ability to recover reliable representations because their denoising process is still constrained by attacked structural evidence.

To address this limitation, we propose DVCL, a Dual-View Contrastive Learning framework for robust heterogeneous graph neural networks. DVCL constructs a purified topology view from meta-path-induced semantic graphs and a feature-induced view from target-node feature similarity. The topology view preserves heterogeneous semantic expressiveness through meta-path modeling, semantic attention, and two-stage edge filtering. The feature-induced view provides a complementary, relatively topology-independent reference under structural attacks. DVCL further aligns the two views with a symmetric cross-view contrastive objective, encouraging each node to maintain consistent representations across semantic topology and feature similarity while remaining discriminative from other nodes.

The main contributions of this paper are summarized as follows:
\begin{itemize}
    \item We propose DVCL, a dual-view robust heterogeneous graph learning framework that jointly exploits purified semantic topology and feature-induced similarity for node classification under structural attacks.
    \item We design a topology purification strategy that combines meta-path-level edge confidence, semantic attention, and global semantic filtering, and pair it with a feature-induced KNN view to reduce reliance on attacked graph structure.
    \item We introduce a symmetric cross-view contrastive objective and a complete evaluation protocol for clean, global-attack, and target-attack settings on heterogeneous graph benchmarks.
\end{itemize}

